\section{至多可数集}

\begin{definition}{至多可数,可数,不可数}{}
	设$E$是集合.若$\left|E\right|<\aleph_{0}$,则称$E$是至多可数集;若$\left|E\right|\leq\aleph_{0}$,则称$E$是可数集;若$\left|E\right|>\aleph_{0}$,则称$E$是不可数集.
\end{definition}

\begin{proposition}{}{}
	\begin{enumerate}
		\item 一个可数集的每个子集都是至多可数集.
		\item 由可数集构成的有限族的乘积是至多可数集.
		\item 由可数集构成的序列$\left(X_{n}\right)_{n=0}^{\infty}$的并是至多可数集.
	\end{enumerate}
\end{proposition}

\begin{proposition}{}{}
	每个至多可数的无限集都与$\N$等势.
\end{proposition}

\begin{proposition}{}{}
	每个无限集$E$都有一个由至多可数的无限子集构成的划分$\left(X_{i}\right)_{i\in I}$,其中$I$和$E$等势.
\end{proposition}

\begin{proposition}{}{}
	设$E$是集合,$F$是无限集,$f:E\to F$是满射.若对任意$y\in F$,$f^{-1}\left(\left\{y\right\}\right)$至多可数,则$F$和$E$等势.
\end{proposition}

\begin{proposition}{}{}
	设$E$是无限集,则$\fin\left(E\right)$和$E$等势.
\end{proposition}

\begin{definition}{连续统的势}{}
	如果一个集合与$\mathcal{P}\left(\N\right)$等势,则称该集合具有连续统的势.
\end{definition}

\begin{remark}{}{}
	显然,具有连续统的势的集合是不可数的.``连续统的势''一词源于$\R$与$\mathcal{P}\left(\N\right)$等势的事实.连续统假设断言,每个不可数集都有一个子集具有连续统的势;广义连续统假设断言,对每个无限基数$\mathfrak{a}$,每个$>\mathfrak{a}$的基数都$\geq 2^{\mathfrak{a}}$.
\end{remark}



























